\documentclass[11pt]{article}

\usepackage[margin=1.5in]{geometry}

\usepackage{amsmath}
\usepackage{lecnotes}
\usetikzlibrary{automata,shapes,positioning,matrix,shapes.callouts,decorations.text,patterns,decorations.pathreplacing,matrix,arrows,chains,calc,tikzmark,petri,topaths}
\pgfdeclarelayer{bg}
\pgfsetlayers{bg,main}

\usepackage{enumerate}

\input{lmacros}

\newcommand{\course}{15-316: Software Foundations of Security \& Privacy}
\newcommand{\lecturer}{Matt Fredrikson}
\newcommand{\lecdate}{April 2, 2026}
\newcommand{\lecnum}{20}
\newcommand{\lectitle}{Trusted Computing}
\newcommand{\courseurl}{https://15316-cmu.github.io/2026/}

% Keys and cryptographic operations
\newcommand{\pk}[1]{\ensuremath{\mi{pk/#1}}}
\newcommand{\sk}[1]{\ensuremath{\mi{sk/#1}}}
\newcommand{\sign}[1]{\ms{sign}_{#1}}
\newcommand{\verify}[1]{\ms{verify}_{#1}}
\newcommand{\cert}[2]{\ms{cert}_{{#1} \arrow {#2}}}
\newcommand{\enc}[1]{\ms{enc}_{{#1}}}

\newcommand{\iskey}[2]{\ms{isKey}(#1,\, #2)}
\newcommand{\pcr}[1]{\ms{PCR}_{#1}}
\newcommand{\extend}[2]{\ms{extend}(#1,\, #2)}
\newcommand{\seal}[3]{\ms{seal}(#1,\, #2,\, #3)}
\newcommand{\readf}[1]{\ms{read}(#1)}

\begin{document}

\maketitle

\noindent\textbf{Sample Solutions for Practice Problems}

\section*{Hash Chains and Authenticated Boot [3 parts]}

We are given that $\pcr{0}$ is initialized to $\mathbf{0}$ and
extended with $D_F$, $D_{BL}$, and $D_{OS}$ in sequence during boot.

\begin{task}\rm
  Write an expression for the final value of $\pcr{0}$.  Explain
  verification and why forgery is infeasible.
\end{task}

\noindent\textbf{Solution.}
After the three extensions:
\begin{align*}
  \text{After } \extend{0}{D_F}: \quad
    & \pcr{0} = H(\mathbf{0} \| D_F) \\
  \text{After } \extend{0}{D_{BL}}: \quad
    & \pcr{0} = H(H(\mathbf{0} \| D_F) \| D_{BL}) \\
  \text{After } \extend{0}{D_{OS}}: \quad
    & \pcr{0} = H(H(H(\mathbf{0} \| D_F) \| D_{BL}) \| D_{OS})
\end{align*}
The remote verifier knows the expected values $D_F$, $D_{BL}$, $D_{OS}$.
It recomputes the hash chain from these values starting from $\mathbf{0}$
and compares the result to the reported $\pcr{0}$.  If they match, the
verifier accepts.

An attacker who wants to claim a different sequence $E_F, E_{BL}, E_{OS}$
(with at least one $E_i \neq D_i$) must find inputs that produce the same
final hash.  This requires finding a second preimage for the composed hash
function: given an input $x$, find $x' \neq x$ such that $H(x') = H(x)$.
Since $H$ is a cryptographic hash function, this is computationally
infeasible.

This relies on the combination of two properties: the extend-only
interface (preventing the attacker from directly writing $\pcr{0}$) and
second preimage resistance (preventing the attacker from finding
alternative inputs that converge to the same hash chain output).

\begin{task}\rm
  An attacker replaces the bootloader with a malicious version that
  measures a fake hash.  Will the verifier accept?
\end{task}

\noindent\textbf{Solution.}
\textbf{No.}  The key insight is that each stage is measured
\emph{before} it runs.  The firmware (which is trusted as the CRTM)
measures the bootloader binary and extends $\pcr{0}$ with
$D'_{BL} = H(\text{malicious bootloader})$ before handing control to it.
Since the malicious bootloader differs from the legitimate one,
$D'_{BL} \neq D_{BL}$.

After the firmware's extend, we have:
\[
  \pcr{0} = H(H(\mathbf{0} \| D_F) \| D'_{BL})
\]
Even though the malicious bootloader then extends with the ``correct''
$D_{OS}$, the final value is:
\[
  \pcr{0} = H(H(H(\mathbf{0} \| D_F) \| D'_{BL}) \| D_{OS})
\]
This differs from the expected value because $D'_{BL} \neq D_{BL}$,
and finding $D'_{BL}$ that produces the same intermediate hash would
require breaking second preimage resistance.

The malicious bootloader can lie about what it measures next, but it
cannot hide the fact that \emph{it itself} was different, because the
previous trusted stage already recorded its hash.

\begin{task}\rm
  The firmware is compromised: it measures the legitimate bootloader
  but hands control to malicious code.  Will the verifier detect this?
\end{task}

\noindent\textbf{Solution.}
\textbf{No.}  The compromised firmware extends $\pcr{0}$ with
$D_F$ for itself, then records $D_{BL}$ for the legitimate bootloader.
But instead of handing control to the legitimate bootloader, it hands
control to malicious code.  The malicious code then extends with
$D_{OS}$ to complete the expected chain.

The final PCR value is:
\[
  \pcr{0} = H(H(H(\mathbf{0} \| D_F) \| D_{BL}) \| D_{OS})
\]
which matches the expected value exactly.  The verifier cannot
distinguish this from a legitimate boot.

This illustrates why the firmware is the root of trust for measurement:
no earlier component measures it, so nothing prevents it from lying
about itself or subsequent stages.  If the CRTM is compromised, the
entire authenticated boot guarantee collapses.  This is precisely why
it must be implemented in tamper-resistant hardware or firmware, and why
it must be as small and simple as possible to minimize the attack
surface.

\clearpage

\section*{Attestation Protocol Design [3 parts]}

\begin{task}\rm
  Describe two distinct attacks against the naive attestation protocol.
\end{task}

\noindent\textbf{Solution.}
\begin{enumerate}
\item \textbf{Replay attack.}  An adversary captures a valid signed
  attestation from a properly configured machine, perhaps by
  eavesdropping on a prior session.  Later, when connecting from a
  compromised machine, the adversary presents the same signed message.
  Since the signature is valid and the PCR value matches the expected
  constant, the server accepts.  The adversary gains access to
  resources intended only for trusted clients.

\item \textbf{Session hijacking (man-in-the-middle).}  Even if the
  attestation is genuine and fresh, the channel between client and
  server is not cryptographically bound to the attestation.  An
  attacker positioned between the two can forward the legitimate
  attestation to the server, causing it to consider the channel
  trusted.  The attacker then injects their own messages into the
  session.  The server has no way to verify that subsequent messages
  come from the same platform that performed the attestation.
\end{enumerate}

\begin{task}\rm
  A nonce is added.  Which attack does it defeat, and which remains?
\end{task}

\noindent\textbf{Solution.}
The nonce defeats the \textbf{replay attack}.  Because $R$ is freshly
generated by the server for each session and is unpredictable, an
attestation signed with a previous nonce will not match the current
one.  The server checks that the hash includes the nonce it issued,
so a captured attestation from a prior session is useless.

The \textbf{session hijacking attack remains}.  After the client sends
the signed attestation (which now includes the fresh nonce), the server
considers the channel trusted.  But nothing binds subsequent messages
to the attested platform.  A man-in-the-middle can relay the nonce to
a legitimate client, collect its fresh attestation, forward it to the
server, and then control the rest of the session.  The nonce proves
the attestation was freshly generated, but does not authenticate
subsequent communication.

Defeating session hijacking requires establishing a shared session key
bound to the attestation, as described in the lecture's full protocol.

\begin{task}\rm
  Explain the privacy concern with using $\sk{tpm}$ directly for two
  unrelated services, and how AIKs address it.
\end{task}

\noindent\textbf{Solution.}
The TPM's endorsement key pair $(\pk{tpm}, \sk{tpm})$ is unique to the
physical hardware and is burned in at manufacture.  If the client uses
$\sk{tpm}$ to sign attestations for both the corporate VPN and the
personal cloud storage provider, then both services see the same
$\pk{tpm}$ in the attestation.  This allows the two services to
collude (or be compromised together) and determine that the same
physical machine---and by extension, likely the same person---is using
both services.  The user's corporate and personal activities become
linkable, violating their privacy.

Attestation identity keys (AIKs) address this by introducing an
indirection layer.  The TPM generates a fresh AIK pair for each
service or context.  The AIK's public key is sent to a Privacy CA,
which verifies that the AIK was generated by a genuine TPM (by
checking the endorsement key certificate).  If valid, the Privacy CA
issues a certificate for the AIK without revealing the endorsement key
to the relying party.

The client uses different AIKs for the VPN and the cloud provider.
Each service sees only the AIK, which is certified as coming from a
genuine TPM but is not linkable to the endorsement key or to AIKs used
with other services.  The Privacy CA must be trusted not to collude
with relying parties, but this single trust point is preferable to
allowing arbitrary cross-service linkage.

This illustrates the tension between accountability (a single global
identity makes attestation simple and strong) and privacy
(unlinkability across contexts protects users from tracking).

\end{document}
